\section{抽象行政行为}
\begin{enumerate}
  \item 记忆为主。考点重复比率高。又叫行政规范性文件。
  \item 行政机关在进行行政管理中，针对不特定的人和事制定普遍适用的规范性文件的活动。
  \item 针对对象是否特定。
  \item 行政法规。
        \begin{enumerate}
          \item 政治类（所有的）、重大/重要的经社文生，及时报告党中央。
          \item 权限。国务院。其他部门发布、经国务院批准也算国务院头上（历史遗留问题），性质属于行政法规。
          \item （立法）程序。从生到死。立、起、审；决、公、生；备、解、评、停。
                \begin{enumerate}
                  \item 立项。国务院下属部门可以，地方不能报请立项。党中央和国务院批准向社会公布。
                  \item 起草。2 个主体。国务院下属部门、法制机构 = 司法部。产物送审稿 \yjt 上下左右。
                        \begin{enumerate}
                          \item 上。拿捏不准分寸，请示国务院批准。
                          \item 下。老百姓、听证会、专家介入、座谈会。\al{可以委托专家、委托组织 \dui}。
                          \item 左右。涉及到其他部门相关领域的，征求意见。
                        \end{enumerate}
                  \item 审查。质检环节。
                        \begin{enumerate}
                          \item 审议说明。法制机构为主，起草部门为辅。\al{审议该规定草案时，须由起草单位作说明 \cuo}。
                          \item 不同意见 \yjt
                                协商 \yjt
                                专家评估 \hz 领导介入 \yjt
                                充分协商 \yjt
                                国务院决定 \hz 协调。
                          \item 送审稿不合格。退回 \hz 缓办。\al{应退回 \cuo。可退回 \dui}！！
                        \end{enumerate}
                  \item 决定。大事开会，小事传批。会 = 常委会。
                  \item 公布。
                        \begin{enumerate}
                          \item 网站、报纸、国务院公报。公报可以有多个，均是标准文本。
                          \item 总理签署、国务院令，军事类加上中央军委主席。
                        \end{enumerate}
                  \item 生效。一般 30 日。金融、外汇、货币政策；国家安全等可立即生效。
                  \item 备案。公布后 30 日内国务院的办公厅报全人常备案。\sjt 日期只有一个数 30。
                  \item 解释。提要求即可，程序很简单。\al{适用行政法规制定程序 \cuo}。
                        \begin{enumerate}
                          \item 补充规定/补充设定。沈报国。进一步明确含义 \hz 出现新情况。省、部级机关报国务院。
                          \item 应用问题。沈法报国。省、部级的法制机构报国务院法制机构。专业对口。
                          \item 解释权属于国务院。
                          \item 行政法规的解释与行政法规具有同等效力。
                        \end{enumerate}
                  \item 评估。国务院部门。
                  \item 停止。国务院决定。\al{政策试点}。
                  \item 冲突。国务院裁决。
                  \item 抵触。\al{某市规章同一行政法规相抵触。向谁提审查建议 \yjt 国务院、省政府 \dui{法制办干活的}}。
                \end{enumerate}
        \end{enumerate}
  \item 地方性法规。重大事项人大过。
        \begin{enumerate}
          \item 协同立法。双方都得地级市以上、门当户对。新增。
          \item 地级市的立法权限。城乡管理、城乡建设、生态文明、历史文化、基层治理。新增。
        \end{enumerate}
  \item 规章。和行政法规高度一致。
        \begin{enumerate}
          \item 重大的经社文生报的是同级党委。\al{杭州市政府 \yjt 杭州市委}。
          \item 制定机关。联合制定 \yjt 首长共同署名公布，最终使用主办机关的命令序号。
          \item 制定程序。决定全体、常务会；评估、法制、上备案；无传批。
                \begin{enumerate}
                  \item 决定。无传批。全体会议 \hz 常务会议。
                  \item 公布。部门公报 \hz 国务院公报、本地/全国的报纸。网络不分本地、全国。
                        \begin{enumerate}
                          \item \al{政府的法制信息网 \dui}。
                          \item \al{该规定须在国务院公报上刊载 \cuo}。
                        \end{enumerate}
                  \item 备案。制定机关的法制机构报上级备案。
                  \item 评估。法制机构负责。没考过。
                  \item 停止。没有。
                \end{enumerate}
          \item 法理内容合并。
                \begin{enumerate}
                  \item 同等效力 \yjt 矛盾统统找国务院。内部省 \dy 市。
                  \item 于公于私俩限制。不能增加权力、减少职责；不能增加义务、减少权利。
                  \item 地级市的规章也是 5 项内容。老人老办法，新人新办法。
                  \item 人大过渡找政府。迫切需要 \yjt 先让政府搞规章 \yjt 2 年为期没有例外 \yjt 法规。同相对保留。
                \end{enumerate}
        \end{enumerate}
\end{enumerate}
